% ============================================================================
% Tabel-tabel untuk BAB 4 (Pengujian & Analisis)
% Cara pakai: % ============================================================================
% Tabel-tabel untuk BAB 4 (Pengujian & Analisis)
% Cara pakai: % ============================================================================
% Tabel-tabel untuk BAB 4 (Pengujian & Analisis)
% Cara pakai: % ============================================================================
% Tabel-tabel untuk BAB 4 (Pengujian & Analisis)
% Cara pakai: \input{tabel_bab4} di bab Anda, atau salin tabel yang diperlukan.
% Sebagian sudah terisi dari hasil notebook; sebagian template (isi sendiri).
% (Tabel pakai tabular standar + \hline; ganti ke booktabs bila template Anda mendukung.)
% ============================================================================

% ---------------------------------------------------------------------------
% 1. ARSITEKTUR & PARAMETER MLP  (SUDAH TERISI)
% ---------------------------------------------------------------------------
\begin{table}[h]
  \centering
  \caption{Arsitektur dan parameter jaringan saraf tiruan (MLP).}
  \label{tab:arsitektur_mlp}
  \begin{tabular}{ll}
    \hline
    \textbf{Aspek} & \textbf{Nilai} \\
    \hline
    Jenis jaringan          & ANN -- Multilayer Perceptron (feedforward) \\
    Arsitektur              & 3--4--1 \\
    Jumlah input            & 3 ($error$, $\Delta error$, $duty$) \\
    Jumlah neuron hidden    & 4 \\
    Jumlah output           & 1 ($\Delta duty$) \\
    Fungsi aktivasi hidden  & tanh \\
    Fungsi aktivasi output  & linear \\
    Normalisasi input       & StandardScaler (mean 0, std 1) \\
    Jumlah parameter        & 21 (16 bobot/bias L1 + 5 L2) \\
    Optimizer pelatihan     & L-BFGS \\
    \hline
  \end{tabular}
\end{table}

% ---------------------------------------------------------------------------
% 2. KOMPOSISI DATASET CLOSED-LOOP  (SUDAH TERISI)
% ---------------------------------------------------------------------------
\begin{table}[h]
  \centering
  \caption{Komposisi dataset pelatihan closed-loop (14 Juni 2026).}
  \label{tab:dataset}
  \begin{tabular}{lc}
    \hline
    \textbf{Keterangan} & \textbf{Jumlah} \\
    \hline
    Jumlah file (sesi)              & 12 \\
    Jumlah episode (global)         & 16 \\
    Baris keputusan ($is\_decision=1$) & 469 \\
    Baris siap latih (X,y)          & 453 \\
    ~~- target naik ($\Delta duty>0$)  & 161 \\
    ~~- target turun ($\Delta duty<0$) & 91 \\
    ~~- target tahan ($\Delta duty=0$) & 201 \\
    Baris setpoint 0.30 bar         & 403 \\
    Baris setpoint 0.25 bar         & 50 \\
    Pembagian Train/Val/Test        & 316 / 46 / 91 \\
    \hline
  \end{tabular}
\end{table}

% ---------------------------------------------------------------------------
% 3. HASIL PELATIHAN MLP (Train/Val/Test)  (SUDAH TERISI)
% ---------------------------------------------------------------------------
\begin{table}[h]
  \centering
  \caption{Hasil pelatihan MLP 3--4--1 pada data latih, validasi, dan uji.}
  \label{tab:hasil_latih}
  \begin{tabular}{lccc}
    \hline
    \textbf{Subset} & \textbf{RMSE (\%)} & \textbf{MAE (\%)} & \textbf{$R^2$} \\
    \hline
    Train      & 0.417 & 0.311 & 0.932 \\
    Validation & 0.541 & 0.334 & 0.882 \\
    Test       & 0.463 & 0.329 & 0.915 \\
    \hline
  \end{tabular}
\end{table}

% ---------------------------------------------------------------------------
% 4. PERBANDINGAN MLP vs BASELINE (TABEL KUNCI)  (SUDAH TERISI)
% ---------------------------------------------------------------------------
\begin{table}[h]
  \centering
  \caption{Perbandingan MLP terhadap baseline pada data uji.}
  \label{tab:baseline}
  \begin{tabular}{lccc}
    \hline
    \textbf{Model} & \textbf{RMSE (\%)} & \textbf{MAE (\%)} & \textbf{$R^2$} \\
    \hline
    MLP 3--4--1                         & \textbf{0.463} & \textbf{0.329} & \textbf{0.915} \\
    Baseline $\Delta duty=0$ (diam)     & 1.597 & 1.121 & $-0.009$ \\
    Baseline Proporsional ($k\cdot e$)  & 1.000 & 0.740 & 0.604 \\
    \hline
  \end{tabular}
\end{table}

% ---------------------------------------------------------------------------
% 5. KARAKTERISASI PLANT: tekanan vs duty per jumlah keran  (TEMPLATE -- ISI)
% ---------------------------------------------------------------------------
\begin{table}[h]
  \centering
  \caption{Karakterisasi plant: tekanan (bar) pada berbagai duty dan jumlah keran.}
  \label{tab:karakterisasi}
  \begin{tabular}{cccccc}
    \hline
    \textbf{Duty (\%)} & \textbf{0 keran} & \textbf{1 keran} & \textbf{2 keran} & \textbf{3 keran} & \textbf{4 keran} \\
    \hline
    70 & \dots & \dots & \dots & \dots & \dots \\
    75 & \dots & \dots & \dots & \dots & \dots \\
    80 & \dots & \dots & \dots & \dots & \dots \\
    85 & \dots & \dots & \dots & \dots & \dots \\
    90 & \dots & \dots & \dots & \dots & \dots \\
    95 & \dots & \dots & \dots & \dots & \dots \\
    \hline
  \end{tabular}
\end{table}

% ---------------------------------------------------------------------------
% 6. PENGUJIAN CLOSED-LOOP DI HARDWARE per skenario  (TEMPLATE -- ISI)
%    settling time, error steady-state, overshoot
% ---------------------------------------------------------------------------
\begin{table}[h]
  \centering
  \caption{Hasil pengujian closed-loop pada perangkat (setpoint 0.30 bar).}
  \label{tab:closedloop}
  \begin{tabular}{cccccc}
    \hline
    \textbf{Keran} & \textbf{Duty awal (\%)} & \textbf{Duty settle (\%)} &
    \textbf{Tekanan settle (bar)} & \textbf{Error steady (bar)} & \textbf{Settling time (s)} \\
    \hline
    4 & \dots & \dots & \dots & \dots & \dots \\
    3 & \dots & \dots & \dots & \dots & \dots \\
    2 & \dots & \dots & \dots & \dots & \dots \\
    1 & \dots & \dots & \dots & \dots & \dots \\
    0 & \dots & \dots & \dots & \dots & \dots \\
    \hline
  \end{tabular}
\end{table}

% ---------------------------------------------------------------------------
% 7. PENGUJIAN SETPOINT TRACKING (berbagai setpoint)  (TEMPLATE -- ISI)
% ---------------------------------------------------------------------------
\begin{table}[h]
  \centering
  \caption{Pengujian penjejakan setpoint (kondisi keran tetap).}
  \label{tab:setpoint_tracking}
  \begin{tabular}{cccc}
    \hline
    \textbf{Setpoint (bar)} & \textbf{Tekanan tercapai (bar)} &
    \textbf{Error (bar)} & \textbf{Keterangan} \\
    \hline
    0.25 & \dots & \dots & \dots \\
    0.30 & \dots & \dots & \dots \\
    0.35 & \dots & \dots & \dots \\
    \hline
  \end{tabular}
\end{table}

% ---------------------------------------------------------------------------
% 8. PENGUJIAN REJEKSI GANGGUAN (perubahan keran saat berjalan)  (TEMPLATE)
% ---------------------------------------------------------------------------
\begin{table}[h]
  \centering
  \caption{Pengujian rejeksi gangguan akibat perubahan jumlah keran (setpoint 0.30 bar).}
  \label{tab:gangguan}
  \begin{tabular}{cccc}
    \hline
    \textbf{Perubahan keran} & \textbf{Deviasi tekanan maks (bar)} &
    \textbf{Waktu pulih (s)} & \textbf{Error akhir (bar)} \\
    \hline
    4 $\rightarrow$ 2 & \dots & \dots & \dots \\
    2 $\rightarrow$ 4 & \dots & \dots & \dots \\
    1 $\rightarrow$ 3 & \dots & \dots & \dots \\
    \hline
  \end{tabular}
\end{table}
 di bab Anda, atau salin tabel yang diperlukan.
% Sebagian sudah terisi dari hasil notebook; sebagian template (isi sendiri).
% (Tabel pakai tabular standar + \hline; ganti ke booktabs bila template Anda mendukung.)
% ============================================================================

% ---------------------------------------------------------------------------
% 1. ARSITEKTUR & PARAMETER MLP  (SUDAH TERISI)
% ---------------------------------------------------------------------------
\begin{table}[h]
  \centering
  \caption{Arsitektur dan parameter jaringan saraf tiruan (MLP).}
  \label{tab:arsitektur_mlp}
  \begin{tabular}{ll}
    \hline
    \textbf{Aspek} & \textbf{Nilai} \\
    \hline
    Jenis jaringan          & ANN -- Multilayer Perceptron (feedforward) \\
    Arsitektur              & 3--4--1 \\
    Jumlah input            & 3 ($error$, $\Delta error$, $duty$) \\
    Jumlah neuron hidden    & 4 \\
    Jumlah output           & 1 ($\Delta duty$) \\
    Fungsi aktivasi hidden  & tanh \\
    Fungsi aktivasi output  & linear \\
    Normalisasi input       & StandardScaler (mean 0, std 1) \\
    Jumlah parameter        & 21 (16 bobot/bias L1 + 5 L2) \\
    Optimizer pelatihan     & L-BFGS \\
    \hline
  \end{tabular}
\end{table}

% ---------------------------------------------------------------------------
% 2. KOMPOSISI DATASET CLOSED-LOOP  (SUDAH TERISI)
% ---------------------------------------------------------------------------
\begin{table}[h]
  \centering
  \caption{Komposisi dataset pelatihan closed-loop (14 Juni 2026).}
  \label{tab:dataset}
  \begin{tabular}{lc}
    \hline
    \textbf{Keterangan} & \textbf{Jumlah} \\
    \hline
    Jumlah file (sesi)              & 12 \\
    Jumlah episode (global)         & 16 \\
    Baris keputusan ($is\_decision=1$) & 469 \\
    Baris siap latih (X,y)          & 453 \\
    ~~- target naik ($\Delta duty>0$)  & 161 \\
    ~~- target turun ($\Delta duty<0$) & 91 \\
    ~~- target tahan ($\Delta duty=0$) & 201 \\
    Baris setpoint 0.30 bar         & 403 \\
    Baris setpoint 0.25 bar         & 50 \\
    Pembagian Train/Val/Test        & 316 / 46 / 91 \\
    \hline
  \end{tabular}
\end{table}

% ---------------------------------------------------------------------------
% 3. HASIL PELATIHAN MLP (Train/Val/Test)  (SUDAH TERISI)
% ---------------------------------------------------------------------------
\begin{table}[h]
  \centering
  \caption{Hasil pelatihan MLP 3--4--1 pada data latih, validasi, dan uji.}
  \label{tab:hasil_latih}
  \begin{tabular}{lccc}
    \hline
    \textbf{Subset} & \textbf{RMSE (\%)} & \textbf{MAE (\%)} & \textbf{$R^2$} \\
    \hline
    Train      & 0.417 & 0.311 & 0.932 \\
    Validation & 0.541 & 0.334 & 0.882 \\
    Test       & 0.463 & 0.329 & 0.915 \\
    \hline
  \end{tabular}
\end{table}

% ---------------------------------------------------------------------------
% 4. PERBANDINGAN MLP vs BASELINE (TABEL KUNCI)  (SUDAH TERISI)
% ---------------------------------------------------------------------------
\begin{table}[h]
  \centering
  \caption{Perbandingan MLP terhadap baseline pada data uji.}
  \label{tab:baseline}
  \begin{tabular}{lccc}
    \hline
    \textbf{Model} & \textbf{RMSE (\%)} & \textbf{MAE (\%)} & \textbf{$R^2$} \\
    \hline
    MLP 3--4--1                         & \textbf{0.463} & \textbf{0.329} & \textbf{0.915} \\
    Baseline $\Delta duty=0$ (diam)     & 1.597 & 1.121 & $-0.009$ \\
    Baseline Proporsional ($k\cdot e$)  & 1.000 & 0.740 & 0.604 \\
    \hline
  \end{tabular}
\end{table}

% ---------------------------------------------------------------------------
% 5. KARAKTERISASI PLANT: tekanan vs duty per jumlah keran  (TEMPLATE -- ISI)
% ---------------------------------------------------------------------------
\begin{table}[h]
  \centering
  \caption{Karakterisasi plant: tekanan (bar) pada berbagai duty dan jumlah keran.}
  \label{tab:karakterisasi}
  \begin{tabular}{cccccc}
    \hline
    \textbf{Duty (\%)} & \textbf{0 keran} & \textbf{1 keran} & \textbf{2 keran} & \textbf{3 keran} & \textbf{4 keran} \\
    \hline
    70 & \dots & \dots & \dots & \dots & \dots \\
    75 & \dots & \dots & \dots & \dots & \dots \\
    80 & \dots & \dots & \dots & \dots & \dots \\
    85 & \dots & \dots & \dots & \dots & \dots \\
    90 & \dots & \dots & \dots & \dots & \dots \\
    95 & \dots & \dots & \dots & \dots & \dots \\
    \hline
  \end{tabular}
\end{table}

% ---------------------------------------------------------------------------
% 6. PENGUJIAN CLOSED-LOOP DI HARDWARE per skenario  (TEMPLATE -- ISI)
%    settling time, error steady-state, overshoot
% ---------------------------------------------------------------------------
\begin{table}[h]
  \centering
  \caption{Hasil pengujian closed-loop pada perangkat (setpoint 0.30 bar).}
  \label{tab:closedloop}
  \begin{tabular}{cccccc}
    \hline
    \textbf{Keran} & \textbf{Duty awal (\%)} & \textbf{Duty settle (\%)} &
    \textbf{Tekanan settle (bar)} & \textbf{Error steady (bar)} & \textbf{Settling time (s)} \\
    \hline
    4 & \dots & \dots & \dots & \dots & \dots \\
    3 & \dots & \dots & \dots & \dots & \dots \\
    2 & \dots & \dots & \dots & \dots & \dots \\
    1 & \dots & \dots & \dots & \dots & \dots \\
    0 & \dots & \dots & \dots & \dots & \dots \\
    \hline
  \end{tabular}
\end{table}

% ---------------------------------------------------------------------------
% 7. PENGUJIAN SETPOINT TRACKING (berbagai setpoint)  (TEMPLATE -- ISI)
% ---------------------------------------------------------------------------
\begin{table}[h]
  \centering
  \caption{Pengujian penjejakan setpoint (kondisi keran tetap).}
  \label{tab:setpoint_tracking}
  \begin{tabular}{cccc}
    \hline
    \textbf{Setpoint (bar)} & \textbf{Tekanan tercapai (bar)} &
    \textbf{Error (bar)} & \textbf{Keterangan} \\
    \hline
    0.25 & \dots & \dots & \dots \\
    0.30 & \dots & \dots & \dots \\
    0.35 & \dots & \dots & \dots \\
    \hline
  \end{tabular}
\end{table}

% ---------------------------------------------------------------------------
% 8. PENGUJIAN REJEKSI GANGGUAN (perubahan keran saat berjalan)  (TEMPLATE)
% ---------------------------------------------------------------------------
\begin{table}[h]
  \centering
  \caption{Pengujian rejeksi gangguan akibat perubahan jumlah keran (setpoint 0.30 bar).}
  \label{tab:gangguan}
  \begin{tabular}{cccc}
    \hline
    \textbf{Perubahan keran} & \textbf{Deviasi tekanan maks (bar)} &
    \textbf{Waktu pulih (s)} & \textbf{Error akhir (bar)} \\
    \hline
    4 $\rightarrow$ 2 & \dots & \dots & \dots \\
    2 $\rightarrow$ 4 & \dots & \dots & \dots \\
    1 $\rightarrow$ 3 & \dots & \dots & \dots \\
    \hline
  \end{tabular}
\end{table}
 di bab Anda, atau salin tabel yang diperlukan.
% Sebagian sudah terisi dari hasil notebook; sebagian template (isi sendiri).
% (Tabel pakai tabular standar + \hline; ganti ke booktabs bila template Anda mendukung.)
% ============================================================================

% ---------------------------------------------------------------------------
% 1. ARSITEKTUR & PARAMETER MLP  (SUDAH TERISI)
% ---------------------------------------------------------------------------
\begin{table}[h]
  \centering
  \caption{Arsitektur dan parameter jaringan saraf tiruan (MLP).}
  \label{tab:arsitektur_mlp}
  \begin{tabular}{ll}
    \hline
    \textbf{Aspek} & \textbf{Nilai} \\
    \hline
    Jenis jaringan          & ANN -- Multilayer Perceptron (feedforward) \\
    Arsitektur              & 3--4--1 \\
    Jumlah input            & 3 ($error$, $\Delta error$, $duty$) \\
    Jumlah neuron hidden    & 4 \\
    Jumlah output           & 1 ($\Delta duty$) \\
    Fungsi aktivasi hidden  & tanh \\
    Fungsi aktivasi output  & linear \\
    Normalisasi input       & StandardScaler (mean 0, std 1) \\
    Jumlah parameter        & 21 (16 bobot/bias L1 + 5 L2) \\
    Optimizer pelatihan     & L-BFGS \\
    \hline
  \end{tabular}
\end{table}

% ---------------------------------------------------------------------------
% 2. KOMPOSISI DATASET CLOSED-LOOP  (SUDAH TERISI)
% ---------------------------------------------------------------------------
\begin{table}[h]
  \centering
  \caption{Komposisi dataset pelatihan closed-loop (14 Juni 2026).}
  \label{tab:dataset}
  \begin{tabular}{lc}
    \hline
    \textbf{Keterangan} & \textbf{Jumlah} \\
    \hline
    Jumlah file (sesi)              & 12 \\
    Jumlah episode (global)         & 16 \\
    Baris keputusan ($is\_decision=1$) & 469 \\
    Baris siap latih (X,y)          & 453 \\
    ~~- target naik ($\Delta duty>0$)  & 161 \\
    ~~- target turun ($\Delta duty<0$) & 91 \\
    ~~- target tahan ($\Delta duty=0$) & 201 \\
    Baris setpoint 0.30 bar         & 403 \\
    Baris setpoint 0.25 bar         & 50 \\
    Pembagian Train/Val/Test        & 316 / 46 / 91 \\
    \hline
  \end{tabular}
\end{table}

% ---------------------------------------------------------------------------
% 3. HASIL PELATIHAN MLP (Train/Val/Test)  (SUDAH TERISI)
% ---------------------------------------------------------------------------
\begin{table}[h]
  \centering
  \caption{Hasil pelatihan MLP 3--4--1 pada data latih, validasi, dan uji.}
  \label{tab:hasil_latih}
  \begin{tabular}{lccc}
    \hline
    \textbf{Subset} & \textbf{RMSE (\%)} & \textbf{MAE (\%)} & \textbf{$R^2$} \\
    \hline
    Train      & 0.417 & 0.311 & 0.932 \\
    Validation & 0.541 & 0.334 & 0.882 \\
    Test       & 0.463 & 0.329 & 0.915 \\
    \hline
  \end{tabular}
\end{table}

% ---------------------------------------------------------------------------
% 4. PERBANDINGAN MLP vs BASELINE (TABEL KUNCI)  (SUDAH TERISI)
% ---------------------------------------------------------------------------
\begin{table}[h]
  \centering
  \caption{Perbandingan MLP terhadap baseline pada data uji.}
  \label{tab:baseline}
  \begin{tabular}{lccc}
    \hline
    \textbf{Model} & \textbf{RMSE (\%)} & \textbf{MAE (\%)} & \textbf{$R^2$} \\
    \hline
    MLP 3--4--1                         & \textbf{0.463} & \textbf{0.329} & \textbf{0.915} \\
    Baseline $\Delta duty=0$ (diam)     & 1.597 & 1.121 & $-0.009$ \\
    Baseline Proporsional ($k\cdot e$)  & 1.000 & 0.740 & 0.604 \\
    \hline
  \end{tabular}
\end{table}

% ---------------------------------------------------------------------------
% 5. KARAKTERISASI PLANT: tekanan vs duty per jumlah keran  (TEMPLATE -- ISI)
% ---------------------------------------------------------------------------
\begin{table}[h]
  \centering
  \caption{Karakterisasi plant: tekanan (bar) pada berbagai duty dan jumlah keran.}
  \label{tab:karakterisasi}
  \begin{tabular}{cccccc}
    \hline
    \textbf{Duty (\%)} & \textbf{0 keran} & \textbf{1 keran} & \textbf{2 keran} & \textbf{3 keran} & \textbf{4 keran} \\
    \hline
    70 & \dots & \dots & \dots & \dots & \dots \\
    75 & \dots & \dots & \dots & \dots & \dots \\
    80 & \dots & \dots & \dots & \dots & \dots \\
    85 & \dots & \dots & \dots & \dots & \dots \\
    90 & \dots & \dots & \dots & \dots & \dots \\
    95 & \dots & \dots & \dots & \dots & \dots \\
    \hline
  \end{tabular}
\end{table}

% ---------------------------------------------------------------------------
% 6. PENGUJIAN CLOSED-LOOP DI HARDWARE per skenario  (TEMPLATE -- ISI)
%    settling time, error steady-state, overshoot
% ---------------------------------------------------------------------------
\begin{table}[h]
  \centering
  \caption{Hasil pengujian closed-loop pada perangkat (setpoint 0.30 bar).}
  \label{tab:closedloop}
  \begin{tabular}{cccccc}
    \hline
    \textbf{Keran} & \textbf{Duty awal (\%)} & \textbf{Duty settle (\%)} &
    \textbf{Tekanan settle (bar)} & \textbf{Error steady (bar)} & \textbf{Settling time (s)} \\
    \hline
    4 & \dots & \dots & \dots & \dots & \dots \\
    3 & \dots & \dots & \dots & \dots & \dots \\
    2 & \dots & \dots & \dots & \dots & \dots \\
    1 & \dots & \dots & \dots & \dots & \dots \\
    0 & \dots & \dots & \dots & \dots & \dots \\
    \hline
  \end{tabular}
\end{table}

% ---------------------------------------------------------------------------
% 7. PENGUJIAN SETPOINT TRACKING (berbagai setpoint)  (TEMPLATE -- ISI)
% ---------------------------------------------------------------------------
\begin{table}[h]
  \centering
  \caption{Pengujian penjejakan setpoint (kondisi keran tetap).}
  \label{tab:setpoint_tracking}
  \begin{tabular}{cccc}
    \hline
    \textbf{Setpoint (bar)} & \textbf{Tekanan tercapai (bar)} &
    \textbf{Error (bar)} & \textbf{Keterangan} \\
    \hline
    0.25 & \dots & \dots & \dots \\
    0.30 & \dots & \dots & \dots \\
    0.35 & \dots & \dots & \dots \\
    \hline
  \end{tabular}
\end{table}

% ---------------------------------------------------------------------------
% 8. PENGUJIAN REJEKSI GANGGUAN (perubahan keran saat berjalan)  (TEMPLATE)
% ---------------------------------------------------------------------------
\begin{table}[h]
  \centering
  \caption{Pengujian rejeksi gangguan akibat perubahan jumlah keran (setpoint 0.30 bar).}
  \label{tab:gangguan}
  \begin{tabular}{cccc}
    \hline
    \textbf{Perubahan keran} & \textbf{Deviasi tekanan maks (bar)} &
    \textbf{Waktu pulih (s)} & \textbf{Error akhir (bar)} \\
    \hline
    4 $\rightarrow$ 2 & \dots & \dots & \dots \\
    2 $\rightarrow$ 4 & \dots & \dots & \dots \\
    1 $\rightarrow$ 3 & \dots & \dots & \dots \\
    \hline
  \end{tabular}
\end{table}
 di bab Anda, atau salin tabel yang diperlukan.
% Sebagian sudah terisi dari hasil notebook; sebagian template (isi sendiri).
% (Tabel pakai tabular standar + \hline; ganti ke booktabs bila template Anda mendukung.)
% ============================================================================

% ---------------------------------------------------------------------------
% 1. ARSITEKTUR & PARAMETER MLP  (SUDAH TERISI)
% ---------------------------------------------------------------------------
\begin{table}[h]
  \centering
  \caption{Arsitektur dan parameter jaringan saraf tiruan (MLP).}
  \label{tab:arsitektur_mlp}
  \begin{tabular}{ll}
    \hline
    \textbf{Aspek} & \textbf{Nilai} \\
    \hline
    Jenis jaringan          & ANN -- Multilayer Perceptron (feedforward) \\
    Arsitektur              & 3--4--1 \\
    Jumlah input            & 3 ($error$, $\Delta error$, $duty$) \\
    Jumlah neuron hidden    & 4 \\
    Jumlah output           & 1 ($\Delta duty$) \\
    Fungsi aktivasi hidden  & tanh \\
    Fungsi aktivasi output  & linear \\
    Normalisasi input       & StandardScaler (mean 0, std 1) \\
    Jumlah parameter        & 21 (16 bobot/bias L1 + 5 L2) \\
    Optimizer pelatihan     & L-BFGS \\
    \hline
  \end{tabular}
\end{table}

% ---------------------------------------------------------------------------
% 2. KOMPOSISI DATASET CLOSED-LOOP  (SUDAH TERISI)
% ---------------------------------------------------------------------------
\begin{table}[h]
  \centering
  \caption{Komposisi dataset pelatihan closed-loop (14 Juni 2026).}
  \label{tab:dataset}
  \begin{tabular}{lc}
    \hline
    \textbf{Keterangan} & \textbf{Jumlah} \\
    \hline
    Jumlah file (sesi)              & 12 \\
    Jumlah episode (global)         & 16 \\
    Baris keputusan ($is\_decision=1$) & 469 \\
    Baris siap latih (X,y)          & 453 \\
    ~~- target naik ($\Delta duty>0$)  & 161 \\
    ~~- target turun ($\Delta duty<0$) & 91 \\
    ~~- target tahan ($\Delta duty=0$) & 201 \\
    Baris setpoint 0.30 bar         & 403 \\
    Baris setpoint 0.25 bar         & 50 \\
    Pembagian Train/Val/Test        & 316 / 46 / 91 \\
    \hline
  \end{tabular}
\end{table}

% ---------------------------------------------------------------------------
% 3. HASIL PELATIHAN MLP (Train/Val/Test)  (SUDAH TERISI)
% ---------------------------------------------------------------------------
\begin{table}[h]
  \centering
  \caption{Hasil pelatihan MLP 3--4--1 pada data latih, validasi, dan uji.}
  \label{tab:hasil_latih}
  \begin{tabular}{lccc}
    \hline
    \textbf{Subset} & \textbf{RMSE (\%)} & \textbf{MAE (\%)} & \textbf{$R^2$} \\
    \hline
    Train      & 0.417 & 0.311 & 0.932 \\
    Validation & 0.541 & 0.334 & 0.882 \\
    Test       & 0.463 & 0.329 & 0.915 \\
    \hline
  \end{tabular}
\end{table}

% ---------------------------------------------------------------------------
% 4. PERBANDINGAN MLP vs BASELINE (TABEL KUNCI)  (SUDAH TERISI)
% ---------------------------------------------------------------------------
\begin{table}[h]
  \centering
  \caption{Perbandingan MLP terhadap baseline pada data uji.}
  \label{tab:baseline}
  \begin{tabular}{lccc}
    \hline
    \textbf{Model} & \textbf{RMSE (\%)} & \textbf{MAE (\%)} & \textbf{$R^2$} \\
    \hline
    MLP 3--4--1                         & \textbf{0.463} & \textbf{0.329} & \textbf{0.915} \\
    Baseline $\Delta duty=0$ (diam)     & 1.597 & 1.121 & $-0.009$ \\
    Baseline Proporsional ($k\cdot e$)  & 1.000 & 0.740 & 0.604 \\
    \hline
  \end{tabular}
\end{table}

% ---------------------------------------------------------------------------
% 5. KARAKTERISASI PLANT: tekanan vs duty per jumlah keran  (TEMPLATE -- ISI)
% ---------------------------------------------------------------------------
\begin{table}[h]
  \centering
  \caption{Karakterisasi plant: tekanan (bar) pada berbagai duty dan jumlah keran.}
  \label{tab:karakterisasi}
  \begin{tabular}{cccccc}
    \hline
    \textbf{Duty (\%)} & \textbf{0 keran} & \textbf{1 keran} & \textbf{2 keran} & \textbf{3 keran} & \textbf{4 keran} \\
    \hline
    70 & \dots & \dots & \dots & \dots & \dots \\
    75 & \dots & \dots & \dots & \dots & \dots \\
    80 & \dots & \dots & \dots & \dots & \dots \\
    85 & \dots & \dots & \dots & \dots & \dots \\
    90 & \dots & \dots & \dots & \dots & \dots \\
    95 & \dots & \dots & \dots & \dots & \dots \\
    \hline
  \end{tabular}
\end{table}

% ---------------------------------------------------------------------------
% 6. PENGUJIAN CLOSED-LOOP DI HARDWARE per skenario  (TEMPLATE -- ISI)
%    settling time, error steady-state, overshoot
% ---------------------------------------------------------------------------
\begin{table}[h]
  \centering
  \caption{Hasil pengujian closed-loop pada perangkat (setpoint 0.30 bar).}
  \label{tab:closedloop}
  \begin{tabular}{cccccc}
    \hline
    \textbf{Keran} & \textbf{Duty awal (\%)} & \textbf{Duty settle (\%)} &
    \textbf{Tekanan settle (bar)} & \textbf{Error steady (bar)} & \textbf{Settling time (s)} \\
    \hline
    4 & \dots & \dots & \dots & \dots & \dots \\
    3 & \dots & \dots & \dots & \dots & \dots \\
    2 & \dots & \dots & \dots & \dots & \dots \\
    1 & \dots & \dots & \dots & \dots & \dots \\
    0 & \dots & \dots & \dots & \dots & \dots \\
    \hline
  \end{tabular}
\end{table}

% ---------------------------------------------------------------------------
% 7. PENGUJIAN SETPOINT TRACKING (berbagai setpoint)  (TEMPLATE -- ISI)
% ---------------------------------------------------------------------------
\begin{table}[h]
  \centering
  \caption{Pengujian penjejakan setpoint (kondisi keran tetap).}
  \label{tab:setpoint_tracking}
  \begin{tabular}{cccc}
    \hline
    \textbf{Setpoint (bar)} & \textbf{Tekanan tercapai (bar)} &
    \textbf{Error (bar)} & \textbf{Keterangan} \\
    \hline
    0.25 & \dots & \dots & \dots \\
    0.30 & \dots & \dots & \dots \\
    0.35 & \dots & \dots & \dots \\
    \hline
  \end{tabular}
\end{table}

% ---------------------------------------------------------------------------
% 8. PENGUJIAN REJEKSI GANGGUAN (perubahan keran saat berjalan)  (TEMPLATE)
% ---------------------------------------------------------------------------
\begin{table}[h]
  \centering
  \caption{Pengujian rejeksi gangguan akibat perubahan jumlah keran (setpoint 0.30 bar).}
  \label{tab:gangguan}
  \begin{tabular}{cccc}
    \hline
    \textbf{Perubahan keran} & \textbf{Deviasi tekanan maks (bar)} &
    \textbf{Waktu pulih (s)} & \textbf{Error akhir (bar)} \\
    \hline
    4 $\rightarrow$ 2 & \dots & \dots & \dots \\
    2 $\rightarrow$ 4 & \dots & \dots & \dots \\
    1 $\rightarrow$ 3 & \dots & \dots & \dots \\
    \hline
  \end{tabular}
\end{table}
